\documentclass[11pt,a4paper]{article}
\usepackage[T1]{fontenc}
\usepackage{lmodern}
\usepackage{amsmath,amssymb}
\usepackage{array,longtable}
\usepackage[margin=25mm]{geometry}
\usepackage{microtype}
\usepackage[hidelinks]{hyperref}
\hypersetup{pdftitle={Literature Review Notes},pdfauthor={Thijs Lenssen}}
\setlength{\parindent}{0pt}
\setlength{\parskip}{0.5em}

\begin{document}
{\LARGE\bfseries Literature Review Notes}

\section*{Paper notes}

Separate reviews are stored in \texttt{Paper\_Notes/}, with an index
in its \texttt{README.md}. Each answers the three per-paper questions;
the synthesis below draws on those reference papers.
Plain-English definitions are collected in the
\hyperref[sec:medical-nomenclature]{medical and neuroscience nomenclature}.

\section*{General questions}

\subsection*{What can or is measured by optogenetics or deep brain stimulation?}
\textbf{Neural activity is measured, and firing rates can be estimated
from those recordings.} In the optogenetic experiments of Zaaimi et al.,
multicontact electrodes record both LFPs and single-neuron spikes;
Fig.~3 reports spike-rate responses during stimulation
\cite[PDF pp.~2, 5]{zaaimi2023}. Electrical-stimulation experiments used
by Nevado Holgado et al. likewise record STN/GPe firing rates, including
GPe responses to STN stimulation \cite[PDF pp.~3--4]{nevado2010}.
The DBS review by Neumann et al. summarizes measured single-neuron
firing responses at several stimulation targets
\cite[PDF pp.~4--6]{neumann2023}.

For a population represented by $n_i$ recorded neurons, a possible
measurement construction is the average spike count per unit time:
\begin{equation*}
\widehat r_i(t)=\frac{1}{n_i\Delta}
\sum_{\ell=1}^{n_i}N_{i\ell}\bigl((t-\Delta,t]\bigr),
\end{equation*}
where $N_{i\ell}$ counts spikes and $\Delta$ is the observation window.
This estimates a population mean when the recorded neurons are
representative. In deviation coordinates it motivates
$y=C_yx+\nu$, with $C_y$ selecting or mixing recorded populations and
$\nu$ representing sampling and measurement error; windowing adds
observation dynamics.

LFPs also provide a route to observing rate dynamics. Zaaimi et al.
model LFP as a high-pass-filtered combination of excitatory and
inhibitory population activity \cite[PDF pp.~9, 15]{zaaimi2023}.
Hall et al., cited by that paper, demonstrate real-time estimation of
individual cortical firing rates from multichannel LFPs using a trained
decoder \cite{hall2014}. Optical indicators offer another route to
activity estimation \cite{grosenick2015,fenno2011}.
Thus both spike-based rate estimates and calibrated indirect
observations can support a firing-rate state model. The observation
map depends on the recording method; an unprocessed LFP is not itself
a firing rate.

\subsection*{Is state feedback sufficient, or is an observer needed?}
\textbf{State feedback is an explicit measurement assumption.}
It is experimentally plausible when the states used by the controller
are population rates estimated from representative, separately
identified recordings of those populations. Sampling error, rate
filtering and latency must then be included in the implementation
model. Recording neural activity alone does not establish access to
every model state.

\textbf{Evidence for estimation.} Zaaimi et al. feed back one filtered
LFP channel, while their model contains both $E$ and $I$
\cite[PDF pp.~2, 9, 14--15]{zaaimi2023}. Their experiment demonstrates
output feedback, rather than separate measurement of both states.
Hall et al. demonstrate an LFP-to-firing-rate decoder
\cite{hall2014}. More directly, Bolus et al. synthesize linear-quadratic
state-feedback gains and implement optogenetic firing-rate regulation
in awake mice using a parameter-adaptive Kalman filter for online
state estimation \cite[Secs.~2.5--2.6]{bolus2021}. This is experimental
precedent for combining state-feedback synthesis with an observer;
their identified linear model differs from the nonlinear delayed model
here.

\textbf{Implication for this project (derived).} For instantaneous
measurements $y=C_yx$, a law $u=Kx$ can be implemented directly if its
required feedback is reconstructible from $y$; for example, if
$K=L C_y$, then $u=Ly$. Otherwise an observer/decoder must estimate
the required state components, provided they can be inferred from the
available measurements and dynamics. A mixed LFP can contain enough
temporal information for estimation, but this requires an observability
or detectability analysis. With an estimate,
\begin{equation*}
u=K\widehat x=Kx+Ke,\qquad e=\widehat x-x.
\end{equation*}
The estimation error enters the closed loop, so the existing
state-feedback certificate alone does not certify the observer-based
nonlinear delayed system. For PIE feedback using delay histories,
those histories are also required; measured rates can be buffered to
supply them after initialization.

\subsection*{What can or is controlled by optogenetics or deep brain stimulation?}
\textbf{Stimulation supplies an input to neuronal populations and changes
their firing dynamics.} Optogenetics targets opsin-expressing neurons
or projections. Zaaimi et al. explicitly represent optical drive $S(t)$
as an input to the excitatory population and measure the resulting
spiking and LFP activity \cite[PDF pp.~2, 5, 15]{zaaimi2023}.
The reviewed optogenetic perturbations of STN/GPe also produce recorded
changes in neuronal firing rates \cite[PDF pp.~6--8]{delacrompe2025}.

For electrical stimulation, Nevado Holgado et al. model STN electrode
current through an additive input and fit its effect to GPe firing-rate
responses \cite[PDF pp.~3--4, Eq.~5]{nevado2010}. DBS recruits neurons
and axons, thereby influencing the stimulated and connected
populations \cite{neumann2023}.
These examples support both input locations in the proposed model:
$c^u u$ changes the drive entering a population's nonlinear activation,
whereas $b^u u$ represents an additive effective rate drive.
Targeting determines which channels are actuated; kinetics and
pathway recruitment determine their gains, signs and dynamics.

\subsection*{What causes the neurological symptoms?}
For Parkinson's disease, dopamine loss changes population firing rates,
synaptic interactions and the temporal organization of basal-ganglia
activity. A dynamical firing-rate model can represent both changes in
mean activity and collective oscillations through time-varying
population rates \cite{nevado2010,lindi2024}. Exaggerated beta
synchronization is associated particularly with bradykinesia and
rigidity; its contribution to symptoms must be distinguished from the
mechanism generating it. Motor deficits can precede exaggerated beta
\cite{delacrompe2025,koeglsperger2026}.

The mechanism is not settled by a single loop: delayed STN--GPe models
and models emphasizing pallidostriatal interactions provide different
accounts of synchronization \cite{nevado2010,lindi2024,ortone2023}.
For freezing of gait, cortical PAC is associated with severity, yet DBS
improves freezing even at comparable PAC levels. Thus PAC reduction
alone does not explain the treatment effect \cite{yin2022}.

\subsection*{What is the control objective?}
At the neural level, the objective can be to regulate population
activity towards a desired rate or trajectory, or to attenuate unwanted
oscillatory/bursting dynamics. Grosenick et al. discuss target neural
activity patterns, and Zaaimi et al. demonstrate feedback enhancement
and suppression of oscillatory activity \cite{grosenick2015,zaaimi2023}.
For regulation about a chosen reference rate, a possible design choice is
\begin{equation*}
x=r-r^{\mathrm{ref}},\qquad
z=\begin{bmatrix}W_r x\\W_u u\end{bmatrix},\qquad
C_z=\begin{bmatrix}W_r\\0\end{bmatrix},\quad
D_{zu}=\begin{bmatrix}0\\W_u\end{bmatrix}.
\end{equation*}
Here the weights trade rate deviations against stimulation effort;
the reference and weights are design choices. For clinical use, the
chosen activity target should be linked to an outcome such as improved
movement initiation or reduced freezing \cite{koehler2024,yin2022}.

\subsection*{Can the system be modeled as follows? If so give the parameterization}
\textbf{Yes: key references explicitly use population firing-rate
models with stimulation inputs and neural observations.}
Zaaimi et al. use a two-population Wilson--Cowan model for their
optogenetic experiments; Nevado Holgado et al. use delayed STN--GPe
rate dynamics constrained by recordings. Azizpour Lindi et al. also
use a rate model alongside a spiking network
\cite{zaaimi2023,nevado2010,lindi2024}.
This supports the proposed model class as a population-level
description; the populations and coefficients depend on the preparation.

Let $x(t)\in\mathbb R^{N_x}$, $w(t)\in\mathbb R^{N_w}$, and
$u(t)\in\mathbb R^{N_u}$ denote the state, disturbance, and control.
Let $v_i(t)$ be the input to $f_i$, $z(t)\in\mathbb R^{N_z}$ the
performance output, and $y(t)\in\mathbb R^{N_y}$ the measurement:
\begin{align*}
\dot{x}_i(t)
 &= a_i x_i(t) + f_i\bigl(v_i(t)\bigr)
    + \sum_{j=1}^{N_x} b^{x}_{ij}x_j(t-\hat{\tau}_{ij}) + \sum_{j=1}^{N_w} b^{w}_{ij}w_j(t-\overline{\tau}_{ij})
         + \sum_{j=1}^{N_u} b^{u}_{ij}u_j(t-\tilde{\tau}_{ij}),\\
v_i(t)
 &= \sum_{j=1}^{N_x} c^{x}_{ij}x_j(t-\hat{\tau}_{ij})
    + \sum_{j=1}^{N_w} c^{w}_{ij}w_j(t-\overline{\tau}_{ij})+ \sum_{j=1}^{N_u} c^{u}_{ij}u_j(t-\tilde{\tau}_{ij}),\\
z(t) &= C_zx(t) + D_{zu}u(t),\\
y_i(t) &= \sum_{j=1}^{N_x} c^{y}_{ij}x_j(t-\tau^{y}_{ij})
          + \sum_{j=1}^{N_w} d^{yw}_{ij}w_j(t-\tau^{yw}_{ij}).
\end{align*}
The first two equations use $i=1,\ldots,N_x$; the measurement equation
uses $i=1,\ldots,N_y$. Local state measurements are obtained by choosing
the corresponding $c^y_{ij}$ nonzero. All delays are constant and
nonnegative, with the required state and input histories prescribed.

\textbf{Nonlinearity assumptions.}
In equilibrium-deviation coordinates, assume $f_i$ is memoryless,
time-invariant, and satisfies $f_i(0)=0$. It is monotone nondecreasing
and slope-restricted on an interval $\mathcal I_i^{\mathrm{sl}}$ containing zero:
\begin{equation*}
0\leq\frac{f_i(v)-f_i(q)}{v-q}\leq\ell_i,
\qquad v,q\in\mathcal I_i^{\mathrm{sl}},\quad v\neq q,\quad 0\leq\ell_i<\infty.
\end{equation*}
On a separate interval $\mathcal I_i^{\mathrm{sec}}$ containing zero,
assume the sector bound
\begin{equation*}
0\leq\frac{f_i(v)}{v}\leq\alpha_i,
\qquad v\in\mathcal I_i^{\mathrm{sec}}\setminus\{0\},\quad 0\leq\alpha_i<\infty.
\end{equation*}
Oddness is not required.

\subsubsection*{Optogenetic example: Zaaimi et al.}
Their Wilson--Cowan model has excitatory and inhibitory activities
$E,I$, optical input $S$, tonic drives $P,Q$ and fluctuating inputs
$\eta_E,\eta_I$ \cite[PDF p.~15, Eqs.~3--4]{zaaimi2023}:
\begin{equation*}
\begin{aligned}
\dot E&=\frac{-E+\sigma(aE-bI+S+P)}{\tau_E}+\eta_E,\\
\dot I&=\frac{-I+\sigma(cE-dI+Q)}{\tau_I}+\eta_I.
\end{aligned}
\end{equation*}
Around an equilibrium, take $x=(E-E^*,I-I^*)^\top$, $u=S-S^*$,
$a_i=-1/\tau_i$, and
\begin{equation*}
[c^x_{ij}]=\begin{bmatrix}a&-b\\c&-d\end{bmatrix},\qquad
c^u=\begin{bmatrix}1\\0\end{bmatrix},\qquad
b^x=b^u=0,\qquad
f_i(v)=\frac{\sigma(q_i^*+v)-\sigma(q_i^*)}{\tau_i}.
\end{equation*}
All intrinsic delays are zero in this example. Additive disturbance
channels can represent fluctuating inputs; the paper simulates them
stochastically. Its logistic has derivative between zero and $1/4$,
giving $\ell_i=\alpha_i=1/(4\tau_i)$ for the centred nonlinearities.
The four letters $a,b,c,d$ above are that paper's coupling parameters;
the numerical simulation values are referred to its Supplementary
Table~2.

The simulated measurement is a high-pass-filtered combination of
$E$ and $I$, passed through the experimental feedback algorithm
\cite[PDF pp.~9, 15]{zaaimi2023}. Thus the paper connects population
dynamics, an optical input and a neural observation in one closed loop.
The filtered measurement can be represented by augmenting the model
with observation-filter states. The recorded spikes additionally allow
rate estimates for the sampled neurons; the two model population
activities are not separately measured as complete populations.

\subsubsection*{Electrical-stimulation example: Nevado Holgado et al.}
Nevado Holgado et al.
\cite[PDF pp.~2--4, Eqs.~2--5 and Tables~1, 3]{nevado2010}
provide a concrete two-population example. Let
$x=(r_S-r_S^*,r_G-r_G^*)^\top$ denote STN and GPe firing-rate deviations
from an equilibrium, and let $q_i^*$ be its total input. Choose
\begin{equation*}
\begin{aligned}
a_i&=-1/\tau_i,& b^x_{ij}&=b^w_{ij}=0,\\
[c^x_{ij}]&=\begin{bmatrix}0&-w_{GS}\\w_{SG}&-w_{GG}\end{bmatrix},&
[\hat\tau_{ij}]&=\begin{bmatrix}0&T_{GS}\\T_{SG}&T_{GG}\end{bmatrix},\\
f_i(v)&=\frac{F_i(q_i^*+v)-F_i(q_i^*)}{\tau_i},&
F_i(q)&=\frac{M_i}{1+\frac{M_i-B_i}{B_i}e^{-4q/M_i}}.
\end{aligned}
\end{equation*}
The nominal cortical and striatal drives are
$(\mathrm{Ctx},\mathrm{Str})=(27,2)$ spikes/s. If their deviations are
included as $w$, take $[c^w_{ij}]=\operatorname{diag}(w_{CS},-w_{XG})$
with zero input delays. The remaining parameters are
\begin{equation*}
\begin{aligned}
(\tau_S,\tau_G,T_{SG},T_{GS},T_{GG})&=(6,14,6,6,4)\ \mathrm{ms},\\
(M_S,B_S,M_G,B_G)&=(300,17,400,75)\ \mathrm{spikes/s},\\
(w_{SG},w_{GS},w_{GG},w_{CS},w_{XG})
&=\begin{cases}
(19,1.12,6.60,2.42,15.1),&\text{healthy fit},\\
(20,10.7,12.3,9.2,139.4),&\text{diseased fit}.
\end{cases}
\end{aligned}
\end{equation*}
Use consistent time units. The sigmoid satisfies $0\leq F_i'\leq1$;
therefore the centred $f_i$ satisfies the stated assumptions globally
with $\ell_i=\alpha_i=1/\tau_i$ (derived from the reported sigmoid).
For their electrode-pulse model, $u$ is current and enters additively
through $b^u=(\mathrm{Stim}/\tau_S,0)^\top$, with $c^u=0$ and zero
input delay; the healthy-data fit gives
$\mathrm{Stim}=4.6\times10^3\,(\mathrm{spikes/s})/\mathrm{mA}$.
This supplies an experimentally fitted example of the additive
stimulation channel.

For spike-derived rate measurements, $C_y$ can select or average the
recorded populations, with sampling error and observation filtering
included as above. If the recorded signal is LFP, its calibrated
observation model must be supplied separately. Performance weights can
be chosen using the proposed $z$ above. More detailed preparations may
require additional populations, opsin states or synaptic filters
\cite{grosenick2015,lindi2024,ortone2023}.

\bigskip
\begingroup
\small
\setlength{\parskip}{0.25em}
\setlength{\abovedisplayskip}{5pt}
\setlength{\belowdisplayskip}{5pt}
\section*{How can I optimally parameterize a PIE Zames--Falb multiplier?}

\textbf{Scope of the answer.} The neural reference papers do not derive
an optimal PIE multiplier. The family below is a derived candidate
whose weights can shape its frequency response. Optimality would be
relative to a specified robustness/performance objective, polynomial
degree and feasible multiplier class.

This derived construction has scalar input and output and distributed
states. Use polynomial profiles on $s\in[0,1]$, a finite number $n$ of
state components, and a scalar feedthrough decision variable $d$, with
\begin{equation*}
\boxed{d>0,\qquad a_j(s)>0,\qquad\beta_j(s)\geq0
\quad\text{for all }s\in[0,1],\ j=1,\ldots,n.}
\end{equation*}
Continuity on this compact interval gives a uniform lower bound
$a_j(s)\geq a_{\min}>0$ and finite integrals.

\subsection*{PIE realization of $M^{-1}$}
Let $Z=L_2([0,1];\mathbb R^n)$. For $x(t,\cdot)\in Z$ and scalar
$u(t),y(t)$, take
\begin{equation*}
\begin{aligned}
\dot x_j(t,s)&=-a_j(s)x_j(t,s)+\beta_j(s)u(t),\qquad x_j(0,s)=0,\\
y(t)&=d\,u(t)+\int_0^1\sum_{j=1}^n a_j(s)x_j(t,s)\,ds.
\end{aligned}
\end{equation*}
Equivalently, $\mathcal T=I$ and
\begin{equation*}
\begin{aligned}
(\mathcal Ax)(s)&=-\operatorname{diag}(a_1(s),\ldots,a_n(s))x(s),&
(\mathcal Bu)(s)&=\begin{bmatrix}\beta_1(s)&\cdots&\beta_n(s)\end{bmatrix}^{\!\top}u,\\
\mathcal Cx&=\int_0^1\sum_j a_j(s)x_j(s)\,ds,& D&=d.
\end{aligned}
\end{equation*}
With the profiles $a_j(s)$ fixed, $\mathcal A,\mathcal C$ are fixed,
and $\mathcal B,D$ are affine in the weight coefficients and $d$.

\subsection*{Transfer functions and normalization}
Writing $p$ for the Laplace variable, define
\begin{equation*}
G(p)=\int_0^1\sum_{j=1}^n\beta_j(s)\frac{a_j(s)}{p+a_j(s)}\,ds.
\end{equation*}
Then
\begin{equation*}
\boxed{
\begin{aligned}
M^{-1}(p)&=d+G(p),& M(p)&=\frac{1}{d+G(p)},\\[0.4em]
H(p)&=\frac{G(p)}{d+G(p)},& H^{-1}(p)&=1+\frac{d}{G(p)}.
\end{aligned}}
\end{equation*}
Use the scaled Zames--Falb normalization $M=d^{-1}(1-H)$;
$d=1$ recovers $M=1-H$. For $G\not\equiv0$, $H^{-1}$ exists
algebraically and may be improper. Spatially varying profiles can give
nonrational transfers; $M,M^{-1},H\in H_\infty$.

\subsection*{Nonnegative impulse response and $L_1$ condition}
Set $b_j(s)=\sqrt{a_j(s)\beta_j(s)}$ and let $\mathcal A_a$ be
multiplication by $\operatorname{diag}(a_j(s))$ on $Z$.
With $b^*x=\langle b,x\rangle_Z$, define
$\mathcal K=\mathcal A_a+d^{-1}bb^*$.
The inversion identity gives
\begin{equation*}
H(p)=\frac1d b^*(pI+\mathcal K)^{-1}b,\qquad
h(t)=\frac1d\langle b,e^{-\mathcal Kt}b\rangle_Z\geq0.
\end{equation*}
The operator $\mathcal K$ is self-adjoint and bounded below by
$a_{\min}I$, so $h$ decays exponentially. Consequently,
\begin{equation*}
\boxed{S:=\int_0^1\sum_j\beta_j(s)\,ds\geq0,\qquad
\|h\|_{L_1}=H(0)=\frac{S}{d+S}<1.}
\end{equation*}
No upper bound on $S$, including $S<d$, is required. This is the
temporal $L_1$ norm of a scalar impulse response. Background:
\href{https://doi.org/10.1109/MCS.2022.3157117}{Scherer (2022),
\emph{Dissipativity and Integral Quadratic Constraints}}, Lemma~31
and (S18), p.~132. Its normalized hard IQC assumes a convex-subgradient
nonlinearity and zero filter initialization; scaling by $1/d>0$
preserves the inequality.

\subsection*{Implication for optimization}
For fixed positive profiles $a_j(s)$, expand $\beta_j(s)$ in a finite
polynomial basis and impose nonnegativity, for example through an
SOS certificate on $[0,1]$. The weight coefficients and $d$ are then
candidate decision variables. An SDP requires the \emph{complete}
IQC/storage inequality to be affine in all joint decision variables;
affinity of $\mathcal B,D$ alone does not establish this. Jointly varying
$a_j(s)$ and a Lyapunov operator generally introduces products of
unknowns. Alternating between fixed-profile and fixed-Lyapunov searches
is therefore a possible heuristic when those subproblems are convex,
with no guarantee of a global optimum.
\endgroup

% Included by literature_review.tex; edit definitions here.
\clearpage
\pdfbookmark[0]{Medical and neuroscience nomenclature}{medical-nomenclature}
\phantomsection\label{sec:medical-nomenclature}
\section*{Medical and neuroscience nomenclature}
Plain-English explanations of the medical and neuroscience terms used
in this review, the paper notes and the related reading. Entries are
grouped by topic. Frequency bands are approximate: individual papers
may use different boundaries.

\begingroup
\small
\setlength{\LTpre}{5pt}
\setlength{\LTpost}{5pt}
\newcommand{\medicalterm}[2]{\textbf{#1} & #2\\[4pt]}
\newcommand{\medicalgroup}[1]{\multicolumn{2}{@{}l}{\rule{0pt}{2.2em}\large\bfseries #1}\\*[6pt]}
\begin{longtable}{@{}>{\raggedright\arraybackslash}p{0.29\linewidth}>{\raggedright\arraybackslash}p{\dimexpr0.71\linewidth-2\tabcolsep\relax}@{}}
\textbf{Term / abbreviation} & \textbf{Meaning}\\\hline
\endfirsthead
\textbf{Term / abbreviation} & \textbf{Meaning (continued)}\\\hline
\endhead
\endfoot
\endlastfoot

\medicalgroup{Neurons and cellular activity}
\medicalterm{Action potential / spike}{A brief electrical pulse generated by a neuron and carried along its axon.}
\medicalterm{Axon}{The extension of a neuron that carries spikes towards other cells.}
\medicalterm{Burst / bursting}{A cluster of spikes occurring close together in time.}
\medicalterm{Cell body / soma}{The main part of a neuron, containing its cell nucleus.}
\medicalterm{Dendrite}{A branching extension through which a neuron receives many of its inputs.}
\medicalterm{Depolarization}{A change that makes the voltage inside a neuron less negative relative to its surroundings.}
\medicalterm{Excitation / excitatory}{An effect that makes a neuron more likely to fire. An excitatory neuron usually has this effect on its targets.}
\medicalterm{Fast-spiking neuron / FSI}{A neuron capable of sustained rapid firing. FSI means fast-spiking interneuron, commonly an inhibitory cell.}
\medicalterm{Firing}{A neuron producing action potentials.}
\medicalterm{Firing rate}{The number of spikes per unit time, usually expressed as spikes per second. It is estimated over a time window.}
\medicalterm{Hyperpolarization}{A change that makes the voltage inside a neuron more negative relative to its surroundings.}
\medicalterm{Inhibition / inhibitory}{An effect that makes a neuron less likely to fire; it need not stop firing completely.}
\medicalterm{Interneuron}{A neuron whose connections mainly influence a local circuit. Many cortical interneurons are inhibitory.}
\medicalterm{Membrane potential}{The voltage difference between the inside and outside of a cell.}
\medicalterm{Neural dynamics}{How neural activity changes over time through inputs and interactions within a circuit.}
\medicalterm{Neural circuit / network}{Interconnected neurons or populations that influence one another.}
\medicalterm{Neural population / ensemble}{A group of neurons considered together, for example because they share a location, cell type or activity pattern.}
\medicalterm{Neuron}{A nerve cell that receives, processes and transmits signals.}
\medicalterm{Neurophysiology}{The study of how neurons and the nervous system function.}
\medicalterm{Parvalbumin neuron / PV neuron}{A neuron containing the calcium-binding protein parvalbumin. Cortical PV neurons are commonly fast-spiking inhibitory interneurons.}
\medicalterm{Population firing rate}{The average firing rate of a group of neurons. A recorded sample can estimate it if the sample represents the group.}
\medicalterm{Pyramidal neuron}{A common excitatory cortical neuron named after its roughly pyramid-shaped cell body; ``pyramidal'' names the cell type, not an oscillation.}
\medicalterm{Spontaneous / endogenous / native activity}{Activity generated by the neural system itself, rather than directly prescribed by an external stimulus.}
\medicalterm{Spike train}{The sequence of times at which a neuron fires.}
\medicalterm{Tonic / phasic activity}{Tonic activity is relatively sustained; phasic activity consists of brief changes or responses.}

\medicalgroup{Synapses, signalling and adaptation}
\medicalterm{Dopamine}{A chemical messenger that changes how neural circuits operate, including circuits involved in movement, motivation and learning.}
\medicalterm{Dopaminergic}{Producing, releasing, responding to or otherwise involving dopamine.}
\medicalterm{Dopamine depletion / hypodopaminergic}{Reduced dopamine availability or signalling in the relevant circuit.}
\medicalterm{D1 / D2 receptors}{Two families of proteins through which dopamine affects cells. Their effects depend on the cell and circuit.}
\medicalterm{GABA}{Gamma-aminobutyric acid: the main inhibitory chemical messenger in the adult brain.}
\medicalterm{GABAergic}{Using GABA to transmit signals.}
\medicalterm{Glutamate / glutamatergic}{Glutamate is the main excitatory chemical messenger in the brain; glutamatergic means using glutamate.}
\medicalterm{Neurotransmitter}{A chemical released by a neuron to communicate with another cell.}
\medicalterm{Postsynaptic potential}{A change in a receiving neuron's membrane voltage caused by synaptic input. It differs from a spike.}
\medicalterm{Presynaptic / postsynaptic}{Presynaptic refers to the sending side of a synapse; postsynaptic refers to the receiving side.}
\medicalterm{Receptor}{A protein that responds to a signalling molecule and changes the receiving cell's activity.}
\medicalterm{Synapse / synaptic}{A connection through which one neuron influences another cell; synaptic means relating to that connection.}
\medicalterm{Synaptic depression}{A decrease in the effect of a synapse during repeated activity. Here ``depression'' does not mean a mood disorder.}
\medicalterm{Synaptic plasticity}{A change in the strength or effectiveness of a synapse over time.}

\medicalgroup{Brain regions and pathways}
\medicalterm{Afferent / efferent}{Afferent connections arrive at a specified region; efferent connections leave it.}
\medicalterm{Arkypallidal neuron / Arky}{A type of GPe neuron with prominent projections back to the striatum.}
\medicalterm{Basal ganglia / BG}{Interconnected brain structures involved in selecting and regulating actions, as well as other functions.}
\medicalterm{Cortex / cortical}{The brain's outer sheet of neural tissue; cortical means relating to it.}
\medicalterm{Cortico-subthalamic}{Connecting or involving the cortex and the subthalamic nucleus.}
\medicalterm{Contralateral / ipsilateral}{On the opposite side / the same side of the body relative to the structure or intervention being described.}
\medicalterm{D1-MSN / D2-MSN}{Striatal neurons distinguished by dopamine-receptor expression, typically associated with the direct and indirect pathways respectively.}
\medicalterm{Direct pathway}{A basal-ganglia route from the striatum to output nuclei such as GPi/SNr, without passing through GPe.}
\medicalterm{Dorsal / ventral}{In the spinal cord, towards the back / the belly side. Dorsoventral means spanning this direction.}
\medicalterm{Fibre / tract / projection}{An axon, bundle of axons or connection carrying signals between neurons or regions.}
\medicalterm{GPe}{External globus pallidus: a basal-ganglia region interacting with the STN, striatum and other regions.}
\medicalterm{GPi}{Internal globus pallidus: a major basal-ganglia output region and a DBS target. It differs from GPe.}
\medicalterm{Hippocampus}{A brain structure important for memory and spatial navigation.}
\medicalterm{Hyperdirect pathway}{A route carrying signals directly from cortex to STN, bypassing the striatum.}
\medicalterm{Indirect pathway}{A basal-ganglia route involving striatum, GPe and downstream output circuits; the classical route also passes through STN.}
\medicalterm{Motor cortex / M1}{Cortical regions involved in movement; M1 is the primary motor cortex.}
\medicalterm{MSN / SPN}{Medium spiny neuron / spiny projection neuron: a principal inhibitory projection neuron of the striatum.}
\medicalterm{Nucleus}{In brain anatomy, a group of neuron cell bodies. This differs from the nucleus inside an individual cell.}
\medicalterm{Pallidal / pallidostriatal}{Pallidal means relating to the globus pallidus; pallidostriatal refers to connections from pallidum to striatum.}
\medicalterm{Premotor cortex / SMA}{Cortical areas involved in preparing and organizing movements; SMA is the supplementary motor area.}
\medicalterm{Prototypic neuron / Proto}{A major GPe cell type that projects to other basal-ganglia regions, including STN.}
\medicalterm{Rostral / caudal}{Towards the head / towards the tail end of the body or spinal cord.}
\medicalterm{Sensorimotor}{Relating to sensory information and movement control.}
\medicalterm{SNc}{Substantia nigra pars compacta: a region containing dopamine-producing neurons that are particularly affected in Parkinson's disease.}
\medicalterm{SNr}{Substantia nigra pars reticulata: a major basal-ganglia output region, distinct from SNc.}
\medicalterm{Spinal cord}{Neural tissue connecting the brain with the body and containing circuits involved in movement and sensation.}
\medicalterm{STN}{Subthalamic nucleus: a small, mainly excitatory basal-ganglia structure and a common DBS target.}
\medicalterm{Striatum / striatal}{A major input region of the basal ganglia; striatal means relating to it.}
\medicalterm{Subcortical}{Located beneath the cerebral cortex.}
\medicalterm{Thalamus / thalamocortical}{The thalamus is a group of brain nuclei communicating with cortex and other regions; thalamocortical means connecting thalamus and cortex.}
\medicalterm{Vim}{Ventral intermediate nucleus of the thalamus: a stimulation target used in tremor treatment.}

\medicalgroup{Neural signals and rhythms}
\medicalterm{Alpha / theta activity}{Named rhythm bands: alpha is roughly 8--12 Hz; theta is roughly 4--8 Hz, extended to 10 Hz in some reviewed work.}
\medicalterm{Beta activity}{Rhythmic activity at roughly 13--30 Hz; some reviewed papers extend the band to 35 Hz.}
\medicalterm{Coherence}{A frequency-specific measure of how consistently two signals are related. It does not by itself establish causation.}
\medicalterm{Entrainment}{A rhythm becoming aligned with an external rhythm, such as periodic stimulation.}
\medicalterm{Evoked response}{Neural activity occurring in response to a stimulus.}
\medicalterm{Gamma activity}{Rhythmic activity above roughly 30 Hz. Upper limits and the definition of high gamma vary by study.}
\medicalterm{Granger causality}{A test of whether one signal's past improves prediction of another. The name does not imply proof of a biological cause.}
\medicalterm{LFP}{Local field potential: a voltage recorded in neural tissue, reflecting combined nearby electrical currents, especially synaptic activity.}
\medicalterm{Neural oscillation / rhythm}{A recurring rise and fall in neural activity. It can occur in population rates or field potentials.}
\medicalterm{PAC}{Phase-amplitude coupling: a statistical relationship between the phase of a slower rhythm and the amplitude of faster activity.}
\medicalterm{Phase / amplitude}{Phase is position within an oscillation's cycle; amplitude describes the size of its variation.}
\medicalterm{Readiness potential}{A slow electrical signal associated with preparation for voluntary movement.}
\medicalterm{Spectral / band power}{How much signal variation lies within a frequency range. Beta power, for example, measures variation in the beta band.}
\medicalterm{Synchronization / synchrony}{Coordination in the timing of neural activity. Neurons need not have identical average firing rates.}
\medicalterm{Wilson--Cowan / firing-rate model}{A population-level description of how neural activity changes through excitation, inhibition and external inputs; Wilson--Cowan models commonly use coupled excitatory and inhibitory populations.}

\medicalgroup{Recording and imaging methods}
\medicalterm{Calcium imaging / indicator}{Recording activity-related changes in intracellular calcium using a fluorescent sensor; calcium signals are slower than individual spikes.}
\medicalterm{Connectome / connectivity}{A map or description of connections between brain regions. Structural connectivity concerns anatomy; functional connectivity concerns statistical relationships between activity signals.}
\medicalterm{ECoG}{Electrocorticography: recording electrical brain activity using electrodes on the cortical surface.}
\medicalterm{EEG}{Electroencephalography: recording electrical brain activity, usually with electrodes on the scalp.}
\medicalterm{Electrode / microelectrode}{A conducting contact used to record or stimulate electrically; a microelectrode is a small contact suited to localized recordings.}
\medicalterm{Electrophysiology}{Measurement of electrical activity in cells or tissue.}
\medicalterm{EMG}{Electromyography: recording electrical activity associated with muscle activation.}
\medicalterm{Extracellular recording}{Recording from outside cells; depending on the setup, this can capture spikes and field potentials.}
\medicalterm{Intracranial recording}{Recording using electrodes placed inside the skull.}
\medicalterm{MEG}{Magnetoencephalography: measuring magnetic fields produced by brain activity.}
\medicalterm{MRI / fMRI}{Magnetic resonance imaging measures tissue properties; functional MRI uses blood-oxygen-related changes to study brain activity indirectly.}
\medicalterm{Multi-unit / single-unit activity}{Multi-unit activity combines spikes from several neurons; single-unit activity consists of spikes assigned to one neuron.}
\medicalterm{PET}{Positron emission tomography: imaging the distribution of an injected radioactive tracer, for example to study molecular systems.}
\medicalterm{Polysomnography / PSG}{A sleep recording combining signals such as EEG, eye movements, muscle activity and breathing.}
\medicalterm{sEEG}{Stereo-electroencephalography: recording electrical activity with depth electrodes placed within the brain.}
\medicalterm{Spike sorting}{Assigning recorded spikes to putative individual neurons using their waveform and other features.}
\medicalterm{Tractography}{Estimating the courses of nerve-fibre bundles from diffusion MRI data.}

\medicalgroup{Stimulation and optogenetics}
\medicalterm{Adaptive DBS / aDBS}{DBS whose settings change automatically in response to measured signals.}
\medicalterm{Antidromic / orthodromic}{Antidromic spikes travel towards the cell body along an axon; orthodromic spikes travel in its usual output direction.}
\medicalterm{Archaerhodopsin / Arch / ArchT}{Light-driven proton pumps used to inhibit neurons; ArchT is an engineered variant.}
\medicalterm{Channelrhodopsin / ChR2 / Chronos}{Light-sensitive ion channels used to excite neurons; ChR2 and Chronos are examples.}
\medicalterm{Closed-loop stimulation}{Stimulation chosen using ongoing measurements. Its purpose may be to increase, suppress or otherwise study activity.}
\medicalterm{DBS}{Deep brain stimulation: electrical stimulation delivered through implanted electrodes to influence brain circuits.}
\medicalterm{Halorhodopsin / NpHR / eNpHR}{A light-driven chloride pump used to inhibit neurons; eNpHR denotes enhanced versions.}
\medicalterm{HFS}{High-frequency stimulation: a rapid train of stimulation pulses; the exact frequency depends on the experiment.}
\medicalterm{Neuromodulation}{Changing neural activity using an intervention such as electrical stimulation, optical stimulation or a chemical agent.}
\medicalterm{Open-loop stimulation}{Stimulation delivered without adjusting it in response to ongoing measurements of the controlled activity.}
\medicalterm{Opsin}{A light-sensitive protein. In these experiments, introduced opsins allow light to influence neuronal activity.}
\medicalterm{Optogenetics}{Using genetically introduced light-sensitive proteins to manipulate selected cells or their connections with light.}
\medicalterm{Optrode}{A device combining optical stimulation with electrode recording.}
\medicalterm{Photostimulation}{Stimulation using light.}
\medicalterm{Projection targeting}{Selecting neurons or their axons according to which region they connect to.}
\medicalterm{Pulse width / stimulation amplitude}{Pulse width is a pulse's duration; amplitude is its electrical or optical strength.}
\medicalterm{Stimulation artifact}{A recorded signal caused by the stimulation or equipment rather than the neural activity being studied.}
\medicalterm{Viral vector / opsin expression}{A modified virus can deliver genetic instructions; opsin expression means cells produce the instructed light-sensitive protein.}

\medicalgroup{Symptoms, disorders and treatment terms}
\medicalterm{Akinesia}{Difficulty initiating movement, or a marked lack of movement.}
\medicalterm{Axial symptoms}{Problems mainly affecting the trunk, posture, balance or walking.}
\medicalterm{Biomarker}{A measurable biological characteristic used to indicate a condition or its response to an intervention.}
\medicalterm{Bradykinesia}{Slowness of movement, often with progressively smaller or slower repeated movements.}
\medicalterm{Cognitive}{Relating to processes such as attention, memory, learning and decision-making.}
\medicalterm{Dyskinesia}{Involuntary movements; in Parkinson's research, often abnormal movements associated with dopaminergic treatment.}
\medicalterm{Dystonia}{Involuntary muscle contractions causing abnormal postures or repetitive twisting movements.}
\medicalterm{Epilepsy / seizure}{A seizure is an episode of abnormal brain activity; epilepsy is a disorder involving an enduring tendency to have seizures.}
\medicalterm{Essential tremor}{A movement disorder characterized mainly by tremor during posture or action.}
\medicalterm{Freezing of gait / FoG}{A brief inability or marked difficulty to step despite the intention to walk.}
\medicalterm{Gait}{The pattern of walking.}
\medicalterm{Huntington's disease}{An inherited neurodegenerative disorder affecting movement, cognition and behaviour.}
\medicalterm{Hypokinetic / hyperkinetic}{Hypokinetic refers to reduced movement; hyperkinetic refers to excessive or unwanted movement.}
\medicalterm{Levodopa / L-DOPA}{A substance converted into dopamine in the body and used to treat Parkinson's symptoms.}
\medicalterm{Motor / non-motor symptoms}{Motor symptoms affect movement; non-motor symptoms include effects on sleep, mood, thinking and other functions.}
\medicalterm{Neurodegeneration}{Progressive damage to or loss of neurons.}
\medicalterm{ON / OFF medication}{Study conditions describing treatment benefit or medication exposure; the paper specifies which meaning is used. These differ from stimulation ON/OFF.}
\medicalterm{Parkinson's disease / PD}{A progressive neurodegenerative disorder involving dopamine-neuron loss and motor as well as non-motor symptoms.}
\medicalterm{Parkinsonism / parkinsonian}{Parkinsonism is a pattern of movement signs resembling Parkinson's disease; parkinsonian also describes these features in experimental models.}
\medicalterm{Pathological / physiological}{Pathological concerns disease-related dysfunction; physiological concerns normal biological function.}
\medicalterm{Phenotype}{An observable characteristic, such as a movement pattern or disease-related behaviour.}
\medicalterm{Rigidity}{Increased resistance when a joint is moved by someone else; commonly experienced as stiffness.}
\medicalterm{Tourette syndrome}{A condition involving recurrent motor and vocal tics.}
\medicalterm{Tremor}{An involuntary rhythmic shaking movement.}
\medicalterm{Volitional / voluntary movement}{Movement initiated intentionally.}

\medicalgroup{Sleep and experimental terminology}
\medicalterm{Acute / chronic}{Acute usually refers to a short-term condition or experiment; chronic refers to a persistent condition or long-term recording/intervention.}
\medicalterm{Anaesthetized}{Under anaesthesia; the drugs can affect neural activity and responses to stimulation.}
\medicalterm{Brain slice}{A thin piece of brain tissue kept viable outside the animal for experiments.}
\medicalterm{Biophysical}{Describing biological activity through physical mechanisms such as membrane currents and ion channels.}
\medicalterm{Causal / correlational evidence}{Causal evidence supports one process influencing another; correlation describes variables changing together.}
\medicalterm{Histology}{Examining tissue structure, usually with microscopy and stains.}
\medicalterm{Epidural}{Located outside the dura mater, the tough outer membrane surrounding the brain and spinal cord.}
\medicalterm{In silico / in vitro / in vivo}{In silico: in a computer model. In vitro: outside the living organism. In vivo: within a living organism.}
\medicalterm{NHP}{Non-human primate, such as a monkey used in research.}
\medicalterm{NREM sleep / N1--N3}{Non-rapid-eye-movement sleep, subdivided into stages N1, N2 and N3.}
\medicalterm{REM sleep}{Rapid-eye-movement sleep: a sleep state with characteristic brain activity, eye movements and usually reduced muscle tone.}
\medicalterm{Sleep decoding}{Estimating sleep state from recorded signals using an algorithm.}
\medicalterm{Transgenic animal}{An animal whose genetic material has been deliberately altered, for example to express an opsin in selected neurons.}
\medicalterm{6-OHDA / MPTP model}{Experimental Parkinsonian models using substances that damage dopamine systems; their effects do not reproduce every feature of human disease.}
\end{longtable}

\textbf{Sources.} Circuit, recording and optogenetic terminology follows
the reviewed papers, particularly
\cite{grosenick2015,fenno2011,neumann2023,delacrompe2025}.
Additional plain-language checks:
\href{https://www.ninds.nih.gov/health-information/disorders/glossary-neurological-terms}{NINDS neurological glossary},
\href{https://www.ninds.nih.gov/current-research/focus-disorders/parkinsons-disease-research/parkinsons-disease-challenges-progress-and-promise}{NINDS Parkinson's disease overview}, and
\href{https://www.ncbi.nlm.nih.gov/books/NBK11117/}{\emph{Neuroscience}: excitatory and inhibitory postsynaptic potentials}.
\endgroup


\newpage
\begingroup
\small
\begin{thebibliography}{99}
\setlength{\itemsep}{0.2em}
\bibitem{grosenick2015}
Grosenick et al. (2015).
\emph{Closed-Loop and Activity-Guided Optogenetic Control}.
\href{https://doi.org/10.1016/j.neuron.2015.03.034}{Neuron}.

\bibitem{fenno2011}
Fenno et al. (2011).
\emph{The Development and Application of Optogenetics}.
\href{https://doi.org/10.1146/annurev-neuro-061010-113817}{Annual Review of Neuroscience}.

\bibitem{hall2014}
Hall, Nazarpour and Jackson (2014).
\emph{Real-time estimation and biofeedback of single-neuron firing rates
using local field potentials}.
\href{https://doi.org/10.1038/ncomms6462}{Nature Communications}.
Further source cited by Zaaimi et al., ref.~13.

\bibitem{bolus2021}
Bolus, Willats, Rozell and Stanley (2021).
\emph{State-space optimal feedback control of optogenetically driven
neural activity}.
\href{https://doi.org/10.1088/1741-2552/abb89c}{Journal of Neural Engineering 18, 036006}.
Additional source on observer-based implementation.

\bibitem{yin2022}
Yin et al. (2022).
\emph{Cortical phase-amplitude coupling is key to the occurrence and
treatment of freezing of gait}.
\href{https://doi.org/10.1093/brain/awac121}{Brain}.

\bibitem{koehler2024}
K{"o}hler et al. (2024).
\emph{Dopamine and deep brain stimulation accelerate the neural dynamics
of volitional action in Parkinson's disease}.
\href{https://doi.org/10.1093/brain/awae219}{Brain}.

\bibitem{yin2024}
Yin et al. (2024).
\emph{Generalized sleep decoding with basal ganglia signals in multiple
movement disorders}.
\href{https://doi.org/10.1038/s41746-024-01115-7}{npj Digital Medicine}.

\bibitem{liu2018}
Liu et al. (2018).
\emph{A Compact Closed-Loop Optogenetics System Based on Artifact-Free
Transparent Graphene Electrodes}.
\href{https://doi.org/10.3389/fnins.2018.00132}{Frontiers in Neuroscience}.

\bibitem{koeglsperger2026}
Koeglsperger et al. (2026).
\emph{Adaptive Deep Brain Stimulation for Parkinson's Disease:
Navigating the Roadblocks to Clinical Implementation}.
\href{https://doi.org/10.1002/mds.70284}{Movement Disorders}.

\bibitem{zaaimi2023}
Zaaimi et al. (2023).
\emph{Closed-loop optogenetic control of the dynamics of neural activity
in non-human primates}.
\href{https://doi.org/10.1038/s41551-022-00945-8}{Nature Biomedical Engineering}.

\bibitem{kathe2022}
Kathe et al. (2022).
\emph{Wireless closed-loop optogenetics across the entire dorsoventral
spinal cord in mice}.
\href{https://doi.org/10.1038/s41587-021-01019-x}{Nature Biotechnology}.

\bibitem{neumann2023}
Neumann et al. (2023).
\emph{Neurophysiological mechanisms of deep brain stimulation across
spatiotemporal resolutions}.
\href{https://doi.org/10.1093/brain/awad239}{Brain}.

\bibitem{rajamani2024}
Rajamani et al. (2024).
\emph{Deep brain stimulation of symptom-specific networks in
Parkinson's disease}.
\href{https://doi.org/10.1038/s41467-024-48731-1}{Nature Communications}.

\bibitem{delacrompe2025}
De la Crompe et al. (2025).
\emph{Reevaluating the role of the STN-GPe network in generating beta
synchronized oscillatory activity in parkinsonism}.
\href{https://doi.org/10.1016/B978-0-443-21992-4.00013-0}{Handbook of Parkinson's Disease Mechanisms, ch.~17}.

\bibitem{nevado2010}
Nevado Holgado et al. (2010).
\emph{Conditions for the Generation of Beta Oscillations in the
Subthalamic Nucleus-Globus Pallidus Network}.
\href{https://doi.org/10.1523/JNEUROSCI.0817-10.2010}{Journal of Neuroscience}.

\bibitem{lindi2024}
Azizpour Lindi et al. (2024).
\emph{Synaptic Changes in Pallidostriatal Circuits Observed in the
Parkinsonian Model Triggers Abnormal Beta Synchrony with Accurate
Spatio-temporal Properties across the Basal Ganglia}.
\href{https://doi.org/10.1523/JNEUROSCI.0419-23.2023}{Journal of Neuroscience}.

\bibitem{ortone2023}
Ortone et al. (2023).
\emph{Dopamine depletion leads to pathological synchronization of
distinct basal ganglia loops in the beta band}.
\href{https://doi.org/10.1371/journal.pcbi.1010645}{PLOS Computational Biology}.
\end{thebibliography}
\endgroup

\end{document}
